% Part: proof-theory
% Chapter: natural-deduction
% Section: quantifiers

\documentclass[../../../include/open-logic-section]{subfiles}

\begin{document}

\olfileid{pt}{nat}{qua}

\olsection{في انتظام \usetoken{P}{proof} والتعويض}

تنشأ من قواعد الأسوار صعوبات سببها «شروط المتغير المميز» في
$\Elim{\lexists}$ و$\Intro{\lforall}$. وأكثرها يزول إذا ضمنّا ألا يعاد
استعمال~!!{constant}~${\OLId{latin-ordinary}{c}}$ في هذه الاستدلالات. ومن أجل ذلك نضع التعريف
الآتي.

\begin{defn}
نسمي~!!{proof} الاستنباط الطبيعي~$\delta$ \emph{منتظمًا} إذا استوفى
الشرطين الآتيين:
\begin{enumerate}
  \item ألا يستعمل متغير مميز~${\OLId{latin-ordinary}{a}}$ بهذا الوصف في أكثر من استدلال واحد
  لـ$\Elim{\lexists}$ أو~$\Intro{\lforall}$؛
  \item وإذا كان~${\OLId{latin-ordinary}{a}}$ هو المتغير المميز في استدلال~$\Elim{\lexists}$
  أو~$\Intro{\lforall}$، ألا يظهر إلا في~!!{proof} الجزئي المباشر المؤدي
  إلى المقدمة الصغرى أو إلى المقدمة، على الترتيب.
\end{enumerate}
\end{defn}

\begin{lem}\ollabel{lem:subst-regular} إذا كان~$\delta$ منتظمًا،
وينتهي بـ$!C$، وكان~${\OLId{latin-ordinary}{c}}$~!!a{constant} لا يُستخدم متغيرًا مميزًا في
$\delta$، وكان~${\OLId{latin-ordinary}{t}}$ حدًا لا يحتوي أي متغير مميز من~$\delta$، فإن
$\Subst{\delta}{{\OLId{latin-ordinary}{t}}}{{\OLId{latin-ordinary}{c}}}$ الناتج من~$\delta$ باستبدال كل ظهور لـ${\OLId{latin-ordinary}{c}}$ بـ
${\OLId{latin-ordinary}{t}}$ هو~!!{proof} منتظم. وتكون~!!{formula} نهاية
$\Subst{\delta}{{\OLId{latin-ordinary}{t}}}{{\OLId{latin-ordinary}{c}}}$ هي~$\Subst{!C}{{\OLId{latin-ordinary}{t}}}{{\OLId{latin-ordinary}{c}}}$. وإذا كان~$!A^{\OLId{latin-ordinary}{x}}$ فرضًا
مفتوحًا في~$\delta$، كان~$\Subst{!A}{{\OLId{latin-ordinary}{t}}}{{\OLId{latin-ordinary}{c}}}^{\OLId{latin-ordinary}{x}}$ فرضًا مفتوحًا في
$\Subst{\delta}{{\OLId{latin-ordinary}{t}}}{{\OLId{latin-ordinary}{c}}}$.
\end{lem}

\begin{proof}
  نسلك الاستقراء على~$\pheight{\delta}$. فإن كان~$\pheight{\delta}={\OLMathNumeral{0}}$
  لم يكن فيه غير الفرض~$!A^{\OLId{latin-ordinary}{x}}$، وعندئذ يكون~$\Subst{\delta}{{\OLId{latin-ordinary}{t}}}{{\OLId{latin-ordinary}{c}}}$ هو
  $\Subst{!A}{{\OLId{latin-ordinary}{t}}}{{\OLId{latin-ordinary}{c}}}^{\OLId{latin-ordinary}{x}}$.

  وإن كان~$\pheight{\delta}>{\OLMathNumeral{0}}$ قسمنا بحسب آخر استدلال في~$\delta$.
  فأما قواعد الروابط القضوية فحالاتها يسيرة نتركها تمارين، وأما هنا
  فنبحث الحالات التي يكون آخر~$\delta$ فيها استدلالًا سُوَريًا.
  \begin{enumerate}
    \item إذا كان آخر الاستدلال~$\Intro{\lforall}$، كان~$\delta$
    على الصورة
    \[
    \AxiomC{}
    \RightLabel{$\delta'$}
    \DeduceC{$!A({\OLId{latin-ordinary}{a}},{\OLId{latin-ordinary}{c}})$}
    \RightLabel{\Intro{\lforall}}
    \UnaryInfC{$\lforall[{\OLId{latin-ordinary}{x}}][!A({\OLId{latin-ordinary}{x}},{\OLId{latin-ordinary}{c}})]$}
    \DisplayProof
    \]
    وبفرض الاستقراء يكون~$\Subst{\delta'}{{\OLId{latin-ordinary}{t}}}{{\OLId{latin-ordinary}{c}}}$~!!{proof} منتظمًا لـ
    $!A({\OLId{latin-ordinary}{a}},{\OLId{latin-ordinary}{t}})$. ولأن~${\OLId{latin-ordinary}{a}}$ متغير مميز في~$\delta$، فلا تظهر~${\OLId{latin-ordinary}{a}}$ في~${\OLId{latin-ordinary}{t}}$.
    ومن ثم يكون آخر الاستدلال في~$\Subst{\delta}{{\OLId{latin-ordinary}{t}}}{{\OLId{latin-ordinary}{c}}}$،
    \[
    \AxiomC{}
    \RightLabel{$\Subst{\delta'}{{\OLId{latin-ordinary}{t}}}{{\OLId{latin-ordinary}{c}}}$}
    \DeduceC{$!A({\OLId{latin-ordinary}{a}},{\OLId{latin-ordinary}{t}})$}
    \RightLabel{\Intro{\lforall}}
    \UnaryInfC{$\lforall[{\OLId{latin-ordinary}{x}}][!A({\OLId{latin-ordinary}{x}},{\OLId{latin-ordinary}{t}})]$}
    \DisplayProof
    \]
    صحيحًا: لما كان~${\OLId{latin-ordinary}{t}}$ لا يحتوي~${\OLId{latin-ordinary}{a}}$، لم تحتوها النتيجة ولا أي فرض
    مفتوح.
    \item إذا كان آخر الاستدلال~$\Elim{\lforall}$، كان~$\delta$
    على الصورة
    \[
    \AxiomC{}
    \RightLabel{$\delta'$}
    \DeduceC{$\lforall[{\OLId{latin-ordinary}{x}}][!A({\OLId{latin-ordinary}{x}},{\OLId{latin-ordinary}{c}})]$}
    \RightLabel{\Elim{\lforall}}
    \UnaryInfC{$!A({\OLId{latin-ordinary}{s}}({\OLId{latin-ordinary}{c}}),{\OLId{latin-ordinary}{c}})$}
    \DisplayProof
    \]
    وبفرض الاستقراء يكون~$\Subst{\delta'}{{\OLId{latin-ordinary}{t}}}{{\OLId{latin-ordinary}{c}}}$~!!{proof} منتظمًا لـ
    $\lforall[{\OLId{latin-ordinary}{x}}][!A({\OLId{latin-ordinary}{x}},{\OLId{latin-ordinary}{t}})]$. (وقد يحتوي الحد~${\OLId{latin-ordinary}{s}}({\OLId{latin-ordinary}{c}})$ في النتيجة~${\OLId{latin-ordinary}{c}}$،
    وقد لا يحتويها.) ويكون آخر الاستدلال في
    $\Subst{\delta}{{\OLId{latin-ordinary}{t}}}{{\OLId{latin-ordinary}{c}}}$،
    \[
    \AxiomC{}
    \RightLabel{$\Subst{\delta'}{{\OLId{latin-ordinary}{t}}}{{\OLId{latin-ordinary}{c}}}$}
    \DeduceC{$\lforall[{\OLId{latin-ordinary}{x}}][!A({\OLId{latin-ordinary}{x}},{\OLId{latin-ordinary}{t}})]$}
    \RightLabel{\Elim{\lforall}}
    \UnaryInfC{$!A({\OLId{latin-ordinary}{s}}({\OLId{latin-ordinary}{t}}),{\OLId{latin-ordinary}{t}})$}
    \DisplayProof
    \]
    صحيحًا، ويكون~$!A({\OLId{latin-ordinary}{s}}({\OLId{latin-ordinary}{t}}),{\OLId{latin-ordinary}{t}}) \ident \Subst{!A({\OLId{latin-ordinary}{s}}({\OLId{latin-ordinary}{c}}),{\OLId{latin-ordinary}{c}})}{{\OLId{latin-ordinary}{t}}}{{\OLId{latin-ordinary}{c}}}$.
    وحالة~$\Intro{\lexists}$ مماثلة.
  \item إذا كان آخر الاستدلال~$\Elim{\lexists}$، كان~$\delta$
    على الصورة
    \[
    \AxiomC{}
    \RightLabel{$\delta_{\OLMathNumeral{1}}'$}
    \DeduceC{$\lexists[{\OLId{latin-ordinary}{x}}][!A({\OLId{latin-ordinary}{x}},{\OLId{latin-ordinary}{c}})]$}
    \AxiomC{$\Discharge{!A({\OLId{latin-ordinary}{a}},{\OLId{latin-ordinary}{c}})}{{\OLId{latin-ordinary}{x}}}$}
    \RightLabel{$\delta_{\OLMathNumeral{2}}'$}
    \DeduceC{$!C$}
    \DischargeRule{\Elim{\lexists}}{{\OLId{latin-ordinary}{x}}}
    \BinaryInfC{$!C$}
    \DisplayProof
    \]
    وبفرض الاستقراء يكون~$\Subst{\delta_{\OLMathNumeral{1}}'}{{\OLId{latin-ordinary}{t}}}{{\OLId{latin-ordinary}{c}}}$ و
    $\Subst{\delta_{\OLMathNumeral{2}}'}{{\OLId{latin-ordinary}{t}}}{{\OLId{latin-ordinary}{c}}}$~!!{proof}s منتظمين: الأول لـ
    $\lexists[{\OLId{latin-ordinary}{x}}][!A({\OLId{latin-ordinary}{x}},{\OLId{latin-ordinary}{t}})]$، والثاني لـ$\Subst{!C}{{\OLId{latin-ordinary}{t}}}{{\OLId{latin-ordinary}{c}}}$ من الفرض
    المفتوح~$!A({\OLId{latin-ordinary}{a}},{\OLId{latin-ordinary}{t}})^{\OLId{latin-ordinary}{x}}$. ولأن~${\OLId{latin-ordinary}{a}}$ متغير مميز في~$\delta$، فلا تظهر
    ${\OLId{latin-ordinary}{a}}$ في~${\OLId{latin-ordinary}{t}}$. ومن ثم يكون آخر الاستدلال في
    $\Subst{\delta}{{\OLId{latin-ordinary}{t}}}{{\OLId{latin-ordinary}{c}}}$،
    \[
    \AxiomC{}
    \RightLabel{$\Subst{\delta_{\OLMathNumeral{1}}'}{{\OLId{latin-ordinary}{t}}}{{\OLId{latin-ordinary}{c}}}$}
    \DeduceC{$\lexists[{\OLId{latin-ordinary}{x}}][!A({\OLId{latin-ordinary}{x}},{\OLId{latin-ordinary}{t}})]$}
    \AxiomC{$\Discharge{!A({\OLId{latin-ordinary}{a}},{\OLId{latin-ordinary}{t}})}{{\OLId{latin-ordinary}{x}}}$}
    \RightLabel{$\Subst{\delta_{\OLMathNumeral{2}}'}{{\OLId{latin-ordinary}{t}}}{{\OLId{latin-ordinary}{c}}}$}
    \DeduceC{$\Subst{!C}{{\OLId{latin-ordinary}{t}}}{{\OLId{latin-ordinary}{c}}}$}
    \DischargeRule{\Elim{\lexists}}{{\OLId{latin-ordinary}{x}}}
    \BinaryInfC{$\Subst{!C}{{\OLId{latin-ordinary}{t}}}{{\OLId{latin-ordinary}{c}}}$}
    \DisplayProof
    \]
    صحيحًا: لا تظهر~${\OLId{latin-ordinary}{a}}$ في المقدمة الكبرى
    $\lexists[{\OLId{latin-ordinary}{x}}][!A({\OLId{latin-ordinary}{x}},{\OLId{latin-ordinary}{t}})]$، ولا في النتيجة~$\Subst{!C}{{\OLId{latin-ordinary}{t}}}{{\OLId{latin-ordinary}{c}}}$، ولا في
    أي فرض مفتوح سوى الفرض المبرأ~$!A({\OLId{latin-ordinary}{a}},{\OLId{latin-ordinary}{t}})^{\OLId{latin-ordinary}{x}}$.
  \end{enumerate}
\end{proof}

\begin{prop}\ollabel{prop:regular}
إذا كان~$\delta$~!!a{proof} في~\Log{N1c} أو~\Log{N1i}، وجد
!!{proof} منتظم~$\delta'$ مماثل له في البنية، ومن ثم في~!!{height}، ولا
يفارق~$\delta$ إلا في إبدال المتغيرات المميزة.
\end{prop}

\begin{proof}
نسمي، لهذا البرهان وحده، استدلال المتغير المميز في~$\delta$
\emph{نظيفًا} إذا لم يستعمل متغيره المميز في استدلال آخر بهذا الوصف،
ولم يظهر ذلك المتغير في~$\delta$ خارج~!!{proof} الاستدلال الجزئي
المباشر؛ وما عدا ذلك فنسميه \emph{متسخًا}. وليكن~${\OLId{latin-ordinary}{n}}$ عدد الاستدلالات
المتسخة في~$\delta$. ونبرهن بالاستقراء القوي على~${\OLId{latin-ordinary}{n}}$ أن~$\delta$ يحول إلى
!!{proof} المنتظم المطلوب~$\delta'$. فإذا كان~${\OLId{latin-ordinary}{n}}={\OLMathNumeral{0}}$ كانت استدلالات
المتغير المميز في~$\delta$ كلها نظيفة، فكان~$\delta$ منتظمًا وانتهى
الأمر. وإن لم يكن، أخذنا أعلى استدلال متسخ لمتغير مميز، وليكن مثلًا
$\Intro{\lforall}$. وعندئذ يكون~!!{proof} الجزئي~$\delta_{\OLMathNumeral{1}}$ من
$\delta$ المنتهي به على الصورة
    \[
    \AxiomC{}
    \RightLabel{$\delta_{\OLMathNumeral{2}}$}
    \DeduceC{$!A({\OLId{latin-ordinary}{c}})$}
    \RightLabel{\Intro{\lforall}}
    \UnaryInfC{$\lforall[{\OLId{latin-ordinary}{x}}][!A({\OLId{latin-ordinary}{x}})]$}
    \DisplayProof
    \]
وتنبّه إلى أن~${\OLId{latin-ordinary}{c}}$ لا تظهر في~$\lforall[{\OLId{latin-ordinary}{x}}][!A({\OLId{latin-ordinary}{x}})]$ ولا في فرض مفتوح لأنها
متغير مميز. ويكون البرهان~$\delta_{\OLMathNumeral{2}}$ منتظمًا، لأن كل استدلال لمتغير
مميز فيه نظيف بحكم اختيارنا أعلى استدلال متسخ. ولتكن~${\OLId{latin-ordinary}{a}}$~!!a{constant}
لا تظهر في~$\delta$.
وبفضل~\olref{lem:subst-regular} يكون~$\Subst{\delta_{\OLMathNumeral{2}}}{{\OLId{latin-ordinary}{a}}}{{\OLId{latin-ordinary}{c}}}$ برهانًا
منتظمًا لـ$!A({\OLId{latin-ordinary}{a}})$. وبشرط المتغير المميز لا يحتوي أي فرض مفتوح في
$\delta_{\OLMathNumeral{2}}$ الثابت~${\OLId{latin-ordinary}{c}}$، ومن ثم لا يحتوي أي فرض مفتوح في
$\Subst{\delta_{\OLMathNumeral{2}}}{{\OLId{latin-ordinary}{a}}}{{\OLId{latin-ordinary}{c}}}$ الثابت~${\OLId{latin-ordinary}{a}}$. وعلى هذا يمكننا تطبيق
$\Intro{\lforall}$. وإذا استبدلنا بـ$\delta_{\OLMathNumeral{1}}$ في~$\delta$ البرهان
$\delta_{\OLMathNumeral{1}}'$،
    \[
    \AxiomC{}
    \RightLabel{$\Subst{\delta_{\OLMathNumeral{2}}}{{\OLId{latin-ordinary}{a}}}{{\OLId{latin-ordinary}{c}}}$}
    \DeduceC{$!A({\OLId{latin-ordinary}{a}})$}
    \RightLabel{\Intro{\lforall}}
    \UnaryInfC{$\lforall[{\OLId{latin-ordinary}{x}}][!A({\OLId{latin-ordinary}{x}})]$}
    \DisplayProof
    \]
نكون قد أقمنا استدلالًا نظيفًا مقام استدلال متسخ لمتغير مميز، ولم نغير
إلا المتغير المميز~${\OLId{latin-ordinary}{c}}$ إلى~${\OLId{latin-ordinary}{a}}$. فيبقى في البرهان الناتج عدد من
استدلالات المتغير المميز المتسخة أقل من~${\OLId{latin-ordinary}{n}}$، ومن فرض الاستقراء القوي
تلزم النتيجة.
\end{proof}

\begin{prob}
صف الحالة في برهان~\olref{prop:regular} التي يكون فيها أعلى استدلال
لمتغير مميز هو~$\Elim{\lexists}$.
\end{prob}

ولن نذكر القضية المتقدمة كلما استعملناها، بل نفرض فيما يأتي انتظام كل
!!{proof}s التي ننظر فيها؛ فإن لم تكن منتظمة بدّلنا متغيراتها المميزة
حتى تنتظم.

تصدق~\Olref{lem:subst-regular} و\olref{prop:regular} كذلك على
\Log{N2c} و\Log{N2i}. ونترك البرهانين تمرينين.

\end{document}
